% !TeX root = ../main.tex
\section{Plastic history and the effect of pore pressure on flow}
\label{sec:poroplastic-demonstration}

We first evaluate the model at a single material point to show how plastic
volume change affects the Biot coefficient.  These calculations isolate the
constitutive response from spatial variations and fluid transport.  We prescribe homogeneous deformation
and pore pressure and solve the constitutive update at each increment.
The skeleton and mineral moduli are those listed for the coupled benchmark
in \cref{tab:mandel-parameters}.  The reference solid volume fraction is
\(\phi_{s0}=0.8\), which keeps the prescribed compression path within
\(0<\phi_s<1\); the plastic parameters are
\(M=0.6\), \(\beta=0.4\), \(c_0=20\ \mathrm{MPa}\), and
\(H=100\ \mathrm{MPa}\), matching the yield and hardening law used in the
coupled calculation. These values define a synthetic
geomaterial that dilates during plastic flow and is used for verification.
The coupled elastic energy supplies the stress at each current plastic
state.  Plastic distention changes the mineral volume and hence both the
Biot coefficient and the mineral-pressure contribution to stress, extending
the reversible elastic coupling to a response that depends on loading history.

\subsection{Loading, unloading, and reloading}

The drained path uses \(\mbf F=\operatorname{diag}(1,F_{yy},1)\).
The axial stretch decreases from unity to \(0.8\), increases to \(0.9\),
and then decreases to \(0.76\).  Each segment occupies one unit of the load
parameter.  The constitutive model is rate independent, so this parameter
only orders the imposed states.

\Cref{fig:implicit-poroplastic-history} shows the resulting Biot coefficient
and plastic distention. Loading is initially elastic; yielding activates
dilation and changes the elastic
volume \(J^e=J/a^p\).  After yielding, \(a^p>1\) indicates a plastic
increase in pore volume even though the material point undergoes compressive
loading.  Elastic compression accommodates this dilative plastic contribution
while the total volume remains below its initial value, \(J=F_{yy}<1\).
On unloading, the plastic factor is fixed while the
mineral and skeleton respond elastically.  The Biot coefficient therefore
differs from the virgin elastic value at the same total compression.  Reloading
retraces this elastic unloading path until the previous yield state is reached;
further compression produces additional plastic distention.

\begin{figure}[ht]
 \centering
 \figureinput{implicit_poroplastic_history.pgf}
 \caption{Drained loading, unloading, and reloading of the poroplastic model.  Left: the Biot coefficient plotted against axial compression; the
 gray dashed curve is the virgin elastic response at the same total deformation.
 Right: accumulated plastic distention along the prescribed path with
 isotropic hardening. The stored
 plastic state remains fixed during elastic unloading and reloading and evolves again after reloading
 reaches yield.}
 \label{fig:implicit-poroplastic-history}
\end{figure}

After yielding, the Biot coefficient remains above the virgin elastic value
at equal total compression.  This agrees with \eqref{eq:poroplastic-biot-correction}:
plastic dilation reduces the elastic mineral volume and increases \(B\).

\subsection{Effect of pore pressure on plastic dilation}
\label{sec:poroplastic-b-feedback}

To examine how pore pressure changes the stress driving plastic flow, we compress the material to \(F_{yy}=0.65\)
while increasing pore pressure linearly from zero to a prescribed final
value between zero and \(1\ \mathrm{GPa}\).  This loading range probes the
synthetic law at large compression and pressure.  In the calculation that
uses the derived constitutive law, stress and the Biot coefficient are
evaluated from the same current plastic state.  We refer to this as the
consistent calculation.  A companion calculation follows the
same deformation and pressure history, initial cohesion, and hardening modulus,
but substitutes the virgin elastic Biot coefficient at the current total volume
and pressure only in the transformation from double-prime to single-prime
stress.  Both calculations evaluate the same double-prime constitutive law,
including its mineral-dependent pressure correction, at their respective
current plastic states.  This comparison measures sensitivity to the
Biot coefficient used in the stress transformation.  The companion is a deliberate
modification of that transformation; its stress is not derived from the same
coupled energy as the consistent calculation.
In the consistent law, the two pressure contributions combine into
\eqref{eq:single-prime-drained-pressure-split}, so \(B\) is not an
independently adjustable multiplier of the plastic driving stress.  The modified calculation changes only one contribution, so this
cancellation no longer holds.  Its purpose is to measure the effect of that
change on predicted plastic flow.  Replacing the reversible Biot coefficient by
the finite multiplier \(\alpha\) would be a different comparison.

\Cref{fig:implicit-poroplastic-feedback} shows that the two plastic responses
coincide at zero pressure and separate as pressure increases.  In this example,
at a common candidate plastic state, the larger current Biot coefficient gives a
smaller isotropic contribution \((1-B)pJ\mbf I\) than the virgin coefficient
in \eqref{eq:single-prime-kirchhoff-stress}.  The consistent transformation
therefore gives greater compressive mean stress at that state.  Over the
complete paths, it produces less plastic dilation than the modified
transformation for the selected material.  All plastic volume change in this
example arises from the dilative part of the nonassociated flow rule.
A cap model that permits plastic pore collapse could produce larger plastic
volume changes and a stronger effect of pressure on plastic flow, depending
on the constitutive
parameters and loading path.

\begin{figure}[ht]
 \centering
 \figureinput{implicit_poroplastic_feedback.pgf}
 \caption{Effect of pore pressure on plastic dilation during simultaneous compression
 and pore-pressure loading to \(F_{yy}=0.65\).  Left: the final Biot coefficient and the virgin elastic
 coefficient at the same total deformation and pressure.  Right: final plastic
 distention using the consistent stress transformation and substituting the
 virgin Biot coefficient only in that transformation.  Each point
 represents a complete constitutive loading path.}
 \label{fig:implicit-poroplastic-feedback}
\end{figure}
\FloatBarrier
